\documentclass{beamer}

% --- THEME AND PACKAGES ---
\usetheme{Berlin} % A clean theme with a sidebar for sections
\usecolortheme{whale} % <-- CHANGED Color theme set to whale

\usepackage[utf8]{inputenc}
\usepackage{graphicx} % For including images
\usepackage{amsmath}  % For math equations
\usepackage{tikz}     % <-- ADDED for precise image placement
\usetikzlibrary{positioning}
\usepackage{algorithm}
\usepackage{algpseudocode}
\usepackage{booktabs}
\usepackage{tabularx}
\usepackage[export]{adjustbox}
% --- Define shorthand commands for bold math vectors ---
\usepackage{amsthm,amsmath,amsfonts,bm,xspace}
\def\vx{{\bm{x}}}
\def\veps{{\bm{\epsilon}}}
\def\vy{{\bm{y}}}

\setbeamertemplate{itemize item}{$\bullet$}
\setbeamertemplate{itemize subitem}{$\bullet$}
% --- Puts 1st level enumerate numbers in a circle ---
% --- NEW CODE FOR CIRCLED NUMBERS (Place AFTER tikz packages) ---
% --- Puts 1st level enumerate numbers in a FILLED circle ---
\setbeamertemplate{enumerate item}{%
	\begin{tikzpicture}[baseline=(a.base)]
		\node[circle,
		fill=blue!80!black,  % <-- CHANGED FROM beamerstructure.fg
		text=white,
		inner sep=1pt,
		minimum size=1.2em,
		anchor=center] (a) {\insertenumlabel};
	\end{tikzpicture}%
}

% --- Puts 2nd level enumerate numbers in a FILLED circle ---
\setbeamertemplate{enumerate subitem}{%
	\begin{tikzpicture}[baseline=(a.base)]
		\node[circle,
		fill=blue!60!black,  % <-- CHANGED FROM theme color mix
		text=white,
		inner sep=1pt,
		minimum size=1.2em,
		anchor=center] (a) {\insertenumlabel};
	\end{tikzpicture}%
}

% --- Puts TOC section numbers in a FILLED circle ---
\setbeamertemplate{section in toc}{%
	\leavevmode
	\begin{tikzpicture}[baseline=(a.base)]
		\node[circle,
		fill=blue!80!black,  % Already changed
		text=white,
		inner sep=1pt,
		minimum size=1.2em,
		anchor=center] (a) {\inserttocsectionnumber};
	\end{tikzpicture}%
	\hspace{1ex}%
	\inserttocsection
}


% This information will be used to create the title page
\title[Next-Generation Agricultural Expert System]{ 	Towards a Next-Generation Agricultural Expert System: A Multimodal Conversational Approach}
\subtitle{}
\author[Shajid Ali Khan]{
	Presented by: \textcolor{blue!80!black}{\textbf{Shajid Ali Khan}}\\
	Supervisor: Dr. Deepak Kumar Singh
}
\institute[IIIT Lucknow]{
	% --- MODIFIED AS REQUESTED ---
	Roll No: MCS25049 \\
	\textcolor{blue!70!white}{Indian Institute of Information Technology Lucknow} % Replace with your institute
	% --- END MODIFICATION ---
}
\date{\today}

% --- LOGO STATEMENT ---
% 1. Add your logo file (e.g., "logo.png") to the same folder as your .tex file.
% 2. Uncomment the line below to add your institute's logo to the title page.
\logo{\includegraphics[height=1cm]{Logo_IIITL.png}} % <-- UNCOMMENT AND EDIT


\begin{document}
	
	%----------------------------------------------------------------------
	% SLIDE 1: TITLE PAGE
	%----------------------------------------------------------------------
	\begin{frame}
		\titlepage % This command renders the title page
		% --- Manually add logo to the top right corner ---
		\begin{tikzpicture}[remember picture, overlay]
			\node[anchor=north east, % Anchor point of the node
			xshift=-1cm, % Horizontal margin from right edge
			yshift=-1cm] % Vertical margin from top edge
			at (current page.south east) % Position relative to page corner
			{
				% --- UNCOMMENT AND EDIT THE LINE BELOW ---
				\includegraphics[height=1cm]{Logo_IIITL.png} % Your logo file
			};
		\end{tikzpicture}
	\end{frame}
	
	%----------------------------------------------------------------------
	% SLIDE 2: TABLE OF CONTENTS
	%----------------------------------------------------------------------
	\begin{frame}
		\frametitle{Table of Contents}
		% This command automatically generates the TOC
		% from the \section commands above.
		\tableofcontents
	\end{frame}
	
	% (This code goes after the TOC's \end{frame})

%----------------------------------------------------------------------
% SLIDE 3: CONTENT PAGE 1 (Introduction) - REVISED
%----------------------------------------------------------------------
\section{Introduction} % <-- Label for TOC navigation

\begin{frame}{Introduction: Next-Generation Agricultural Expert System}
	\footnotesize
	\begin{itemize}
		\item \textbf{The Shift:} Moving from static, rule-based apps to AI-driven conversational agents.
		
		\vspace{0.2cm}
		\item \textbf{The Challenge:} Farming requires synthesizing heterogeneous data (visual anomalies, complex soil chemistry, local policies).
		
		\vspace{0.2cm}
		\item \textbf{Current Limitations:} Generic LLMs hallucinate agronomic advice, lack domain specificity, and fail on low-quality field media.
		
		\vspace{0.2cm}
		\item \textbf{Proposed Solution:} A multimodal framework combining computer vision, dynamic knowledge retrieval, and strict domain constraints.
	\end{itemize}
\end{frame}


\section{Related Work}
\begin{frame}{Related Work -- Vision and Multimodality}
	\footnotesize
	\begin{itemize}
		\item \textbf{Multimodal LLMs (e.g., AgriM-LLM):} Attempt to fuse visual encoders with language models for disease management.
		\begin{itemize}
			\item \textit{Limitation:} Accuracy degrades heavily on the blurry, compressed media typical of real-world farming.
		\end{itemize}
		
		\vspace{0.2cm}
		\item \textbf{Diagnostic CNNs:} Proven highly accurate in identifying nutrient deficiencies via static plant imagery.
		\begin{itemize}
			\item \textit{Limitation:} Over-reliant on high-fidelity, static images rather than dynamic, noisy field video.
		\end{itemize}
	\end{itemize}
\end{frame}

\begin{frame}{Related Work -- Architecture and NLP}
	\footnotesize
	\begin{itemize}
		\item \textbf{Retrieval-Augmented Generation (RAG):} The current standard for grounding LLMs in localized data to prevent hallucinations.
		\begin{itemize}
			\item \textit{Limitation:} Current RAG implementations lack temporal memory; they treat queries in isolation instead of tracking a full crop cycle.
		\end{itemize}
		
		\vspace{0.2cm}
		\item \textbf{Lexical Processing:} Standard NLP pipelines often fail to capture localized agricultural vernacular and context. 
		
		\vspace{0.2cm}
		\item \textbf{The Regulatory Void:} Existing chatbots offer chemical/agronomic advice while operating in a legal vacuum, completely bypassing agricultural laws and compliance frameworks.
	\end{itemize}
\end{frame}


\section{Research Gaps}
\begin{frame}{Identified Research Gaps}

	\footnotesize
	\begin{enumerate}
		\item \textbf{Low-Fidelity Media Handling:} Current models cannot accurately classify anomalies from the low-resolution, noisy images or videos captured by standard smartphones in the field.
		
		\vspace{0.2cm}
		\item \textbf{Contextual Memory Loss:} Existing RAG architectures fail to maintain a continuous, evolving state of a farm over an entire harvesting season.
		
		\vspace{0.2cm}
		\item \textbf{The Compliance Disconnect:} There is an absence of expert systems that cross-reference deep biochemical recommendations with regional agricultural regulations and legal constraints.
	\end{enumerate}
\end{frame}

\section{Objectives}
\begin{frame}{Research Objectives}

	\footnotesize
	\begin{enumerate}
		\item \textbf{Enhance Visual Input:} Integrate Super-Resolution (Image and Video) preprocessing to restore low-fidelity field media prior to diagnostic classification.
		
		\medskip
		\item \textbf{Develop Robust Architecture:} Construct an advanced RAG framework equipped with state-tracking to maintain full-season agronomic context.
		
		\medskip
		\item \textbf{Embed Multidisciplinary Intelligence:} Bridge agronomic advice with legal compliance, ensuring all soil and chemical recommendations adhere strictly to local policies.
	\end{enumerate}
\end{frame}

\section*{References}
\begin{frame}[allowframebreaks]{References}
	\scriptsize
	\begin{thebibliography}{10}
		
		\bibitem{taha2022}
		Taha M. F., Abdalla A., ElMasry G., Gouda M., et al. (2022).
		\newblock Using Deep Convolutional Neural Network for Image-Based Diagnosis of Nutrient Deficiencies in Plants Grown in Aquaponics.
		\newblock \emph{Chemosensors}, 10(2), 45.
		
		\bibitem{ahmedm2026}
		Ahmedm Y., \& Luan R. (2026).
		\newblock AgriM-LLM: An Agriculture-Specific Multimodal Large Language Model for Intelligent Crop Disease and Pest Management.
		\newblock \emph{Preprints.org}.
		
		\bibitem{patne2025}
		Patne A. S., Ranjane S. S., Walke O., \& Owhal R. (2025).
		\newblock AgriSaathi: A Comprehensive Survey on Large Language Models and Multimodal AI for Agricultural Advisory Systems.
		\newblock \emph{International Journal For Multidisciplinary Research (IJFMR)}, 7(6).
		
		\bibitem{yin2025}
		Yin S., Xi Y., Zhang X., Sun C., \& Mao Q. (2025).
		\newblock Foundation Models in Agriculture: A Comprehensive Review.
		\newblock \emph{Agriculture}, 15(8).
		
		\bibitem{akbar2025}
		Akbar N. A., Dembani R., Lenzitti B., \& Tegolo D. (2025).
		\newblock RAG-Driven Memory Architectures in Conversational LLMs—A Literature Review With Insights Into Emerging Agriculture Data Sharing.
		\newblock \emph{IEEE Access}, 13.
		
		\bibitem{nanthakumar2025}
		Nanthakumar G., Atchaya R., Brindha S., Harshaa I., \& Kaviya R. (2025).
		\newblock AgriPath: An AI-Driven Assistant Empowering Farmers with Real-Time Guidance, Multilingual Access, and Tailored Farming Support.
		\newblock \emph{International Journal of Advanced Research in Science, Communication and Technology (IJARSCT)}, 5(2).
		
		\bibitem{sirapani2026}
		Sirapani D., Revathi J., \& Ramadevi C. (2026).
		\newblock A SHAP - Enhanced Voice-Based Conversational Agent for Agriculture Using BERT.
		\newblock \emph{Journal of Computer Technology \& Applications}, 17(1).
		
		\bibitem{zhao2022}
		Zhao Z., Feng W., Xiao J., et al. (2022).
		\newblock Rapid and Accurate Prediction of Soil Texture Using an Image-Based Deep Learning Autoencoder Convolutional Neural Network Random Forest (DLAC-CNN-RF) Algorithm.
		\newblock \emph{Agronomy}, 12(12), 3063.
		
		\bibitem{kim2023}
		Kim D., Kim T., Jeon J., \& Son Y. (2023).
		\newblock Soil-Surface-Image-Feature-Based Rapid Prediction of Soil Water Content and Bulk Density Using a Deep Neural Network.
		\newblock \emph{Applied Sciences}, 13(7), 4430.
		
		\bibitem{hossain2023}
		Hossain, M. R. H., \& Kabir, M. A. (2023).
		\newblock Machine Learning Techniques for Estimating Soil Moisture from Smartphone Captured Images.
		\newblock \emph{Agriculture}, 13(3), 574.
		
	\end{thebibliography}
\end{frame}

\section*{End}
\begin{frame}
	\begin{center}
		\vspace{1.5cm}
		\Huge \textbf{Thank You!} 
		
		\vspace{1cm}
		\Large \textit{Any Questions?}
	\end{center}
\end{frame}
\end{document} % <-- THIS MUST BE THE VERY LAST LINE